\documentclass[12pt]{article}
\usepackage[UTF8]{inputenc}
\usepackage{thaisupp}
\usepackage{geometry}
\geometry{a4paper, margin=2.5cm}
\begin{document}
\begin{center}
{\LARGE\textbf{ฟิสิกส์ ม.ปลาย — Part 1}}\\[0.3em]
{\Large\textbf{Lesson 8: โมเมนตัมและการชน}}\\[0.5em]
\rule{0.8\linewidth}{0.5pt}
\end{center}
\section*{คำถาม In-Class}
\begin{enumerate}
\item \textbf{วัตถุมวล 3 kg เคลื่อนที่ด้วยความเร็ว 4 m/s มีโมเมนตัมเท่าใด}
\begin{enumerate}
\item ก) 7 kg·m/s
\item ข) 12 kg·m/s
\item ค) 1.33 kg·m/s
\item ง) 0.75 kg·m/s
\end{enumerate}
\item \textbf{วัตถุมวล 5 kg มีโมเมนตัม 20 kg·m/s วัตถุเคลื่อนที่ด้วยความเร็วเท่าใด}
\begin{enumerate}
\item ก) 4 m/s
\item ข) 25 m/s
\item ค) 100 m/s
\item ง) 0.25 m/s
\end{enumerate}
\item \textbf{วัตถุ A มวล 2 kg วิ่งด้วยความเร็ว 6 m/s ชนกับวัตถุ B ที่อยู่กับที่ (มวล 4 kg) แบบไม่ยืดหยุ่น ความเร็วหลังชนเป็นเท่าใด}
\begin{enumerate}
\item ก) 6 m/s
\item ข) 3 m/s
\item ค) 2 m/s
\item ง) 1.5 m/s
\end{enumerate}
\item \textbf{วัตถุ A มวล 3 kg วิ่งด้วยความเร็ว 4 m/s ชนกับวัตถุ B มวล 2 kg ที่อยู่กับที่ แบบยืดหยุ่น วัตถุ B จะมีความเร็วหลังชนเท่าใด}
\begin{enumerate}
\item ก) 2.4 m/s
\item ข) 3.6 m/s
\item ค) 4.8 m/s
\item ง) 6.0 m/s
\end{enumerate}
\item \textbf{ลูกบอลมวล 0.5 kg กระทบผนังด้วยความเร็ว 10 m/s แล้วสะท้อนกลับด้วยความเร็ว 8 m/s impulse ที่ผนังรับมีค่าเท่าใด}
\begin{enumerate}
\item ก) 1 N·s
\item ข) 4.5 N·s
\item ค) 9 N·s
\item ง) 18 N·s
\end{enumerate}
\item \textbf{ถ้าไม่มีแรงภายนอกกระทำต่อระบบ ข้อใดต่อไปนี้ถูกต้อง}
\begin{enumerate}
\item ก) พลังงานจลน์รวมคงที่
\item ข) โมเมนตัมรวมคงที่
\item ค) ความเร็วรวมคงที่
\item ง) มวลรวมคงที่
\end{enumerate}
\item \textbf{รถบรรทุกมวล 2000 kg วิ่งด้วยความเร็ว 20 m/s ชนกับรถเก๋งมวล 1000 kg ที่อยู่กับที่ แล้วเคลื่อนที่ติดไปด้วยกัน ความเร็วหลังชนเป็นเท่าใด}
\begin{enumerate}
\item ก) 6.67 m/s
\item ข) 10 m/s
\item ค) 13.33 m/s
\item ง) 20 m/s
\end{enumerate}
\item \textbf{หน่วยของโมเมนตัมในระบบ SI คือข้อใด}
\begin{enumerate}
\item ก) kg·m/s²
\item ข) kg·m/s
\item ค) kg·m²/s
\item ง) kg·m/s³
\end{enumerate}
\item \textbf{วัตถุ A และ B มีมวลเท่ากัน วัตถุ A วิ่งด้วยความเร็ว 10 m/s ชนกับ B ที่อยู่กับที่แบบยืดหยุ่น หลังชน A จะมีความเร็วเท่าใด}
\begin{enumerate}
\item ก) 0 m/s
\item ข) 5 m/s
\item ค) -5 m/s
\item ง) 10 m/s
\end{enumerate}
\item \textbf{ลูกบอลมวล 0.2 kg ตกกระทบพื้นด้วยความเร็ว 5 m/s แล้วกระเด้งขึ้นด้วยความเร็ว 3 m/s ในเวลา 0.01 s แรงเฉลี่ยที่พื้นกระทำต่อลูกบอลมีค่าเท่าใด}
\begin{enumerate}
\item ก) 40 N
\item ข) 80 N
\item ค) 120 N
\item ง) 160 N
\end{enumerate}
\item \textbf{ข้อใดต่อไปนี้เป็นลักษณะของการชนแบบไม่ยืดหยุ่น}
\begin{enumerate}
\item ก) พลังงานจลน์คงที่
\item ข) โมเมนตัมคงที่
\item ค) ทั้งโมเมนตัมและพลังงานจลน์คงที่
\item ง) ความเร็วคงที่
\end{enumerate}
\item \textbf{ในการชนแบบยืดหยุ่น ข้อใดต่อไปนี้ถูกต้อง}
\begin{enumerate}
\item ก) โมเมนตัมไม่คงที่ พลังงานจลน์คงที่
\item ข) โมเมนตัมคงที่ พลังงานจลน์คงที่
\item ค) โมเมนตันคงที่ พลังงานจลน์ไม่คงที่
\item ง) โมเมนตัมไม่คงที่ พลังงานจลน์ไม่คงที่
\end{enumerate}
\item \textbf{ปืนมวล 4 kg ยิงกระสุนมวล 0.05 kg ออกไปด้วยความเร็ว 400 m/s ปืนจะเคลื่อนที่ถอยหลังด้วยความเร็วเท่าใด}
\begin{enumerate}
\item ก) 2 m/s
\item ข) 5 m/s
\item ค) 8 m/s
\item ง) 10 m/s
\end{enumerate}
\item \textbf{วัตถุ A มวล 2 kg วิ่งไปทางขวาด้วยความเร็ว 3 m/s และวัตถุ B มวล 4 kg วิ่งไปทางซ้ายด้วยความเร็ว 2 m/s โมเมนตัมรวมเป็นเท่าใด}
\begin{enumerate}
\item ก) 2 kg·m/s ไปทางขวา
\item ข) 2 kg·m/s ไปทางซ้าย
\item ค) 14 kg·m/s ไปทางขวา
\item ง) 14 kg·m/s ไปทางซ้าย
\end{enumerate}
\item \textbf{วัตถุ A มวล 1 kg วิ่งด้วยความเร็ว 6 m/s ชนกับ B มวล 3 kg ที่อยู่กับที่แบบยืดหยุ่น ความเร็วของ A และ B หลังชนเป็นเท่าใด}
\begin{enumerate}
\item ก) A = -3 m/s, B = 3 m/s
\item ข) A = -1.5 m/s, B = 1.5 m/s
\item ค) A = -3 m/s, B = 1.5 m/s
\item ง) A = -4.5 m/s, B = 4.5 m/s
\end{enumerate}
\item \textbf{impulse มีหน่วยเหมือนกับข้อใด}
\begin{enumerate}
\item ก) พลังงาน
\item ข) กำลัง
\item ค) โมเมนตัม
\item ง) แรง
\end{enumerate}
\item \textbf{รถ A มวล 1500 kg วิ่งด้วยความเร็ว 20 m/s ชนกับรถ B มวล 1500 kg ที่วิ่งด้วยความเร็ว 10 m/s ในทิศเดียวกัน แล้วติดกันไป ความเร็วหลังชนเป็นเท่าใด}
\begin{enumerate}
\item ก) 10 m/s
\item ข) 15 m/s
\item ค) 20 m/s
\item ง) 30 m/s
\end{enumerate}
\item \textbf{จรวดมวล 100 kg พ่นแก๊สออกไปด้านหลังด้วยความเร็ว 200 m/s จรวดจะมีความเร็วเท่าใด (สมมติเริ่มต้นอยู่กับที่)}
\begin{enumerate}
\item ก) 50 m/s
\item ข) 100 m/s
\item ค) 150 m/s
\item ง) 200 m/s
\end{enumerate}
\item \textbf{ลูกบอลเทนนิสมวล 0.06 kg ถูกตีด้วยแรงเฉลี่ย 300 N เป็นเวลา 0.02 s ความเร็วของลูกบอลหลังถูกตีเป็นเท่าใด (สมมติเริ่มจากหยุดนิ่ง)}
\begin{enumerate}
\item ก) 50 m/s
\item ข) 75 m/s
\item ค) 100 m/s
\item ง) 150 m/s
\end{enumerate}
\item \textbf{วัตถุ A มวล 2 kg วิ่งด้วยความเร็ว 5 m/s ชนกับ B มวล 3 kg ที่วิ่งด้วยความเร็ว 2 m/s ในทิศตรงข้าม แบบไม่ยืดหยุ่น ความเร็วหลังชนเป็นเท่าใด}
\begin{enumerate}
\item ก) 0.4 m/s ไปทางเดิมของ A
\item ข) 0.4 m/s ไปทางเดิมของ B
\item ค) 1.6 m/s ไปทางเดิมของ A
\item ง) 1.6 m/s ไปทางเดิมของ B
\end{enumerate}
\item \textbf{ในการชนแบบยืดหยุ่น ถ้าวัตถุเบามากๆ (m₂ >> m₁) วิ่งเข้าชนวัตถุเบามากๆ ที่อยู่กับที่ วัตถุเบาจะเคลื่อนที่ออกไปด้วยความเร็วประมาณเท่าใด}
\begin{enumerate}
\item ก) เท่ากับความเร็วเดิมของมัน
\item ข) 2 เท่าของความเร็วเดิม
\item ค) ครึ่งหนึ่งของความเร็วเดิม
\item ง) หยุดนิ่ง
\end{enumerate}
\item \textbf{ถ้าแรงลัพธ์ที่กระทำต่อวัตถุเป็นศูนย์ ข้อใดจะไม่เกิดขึ้น}
\begin{enumerate}
\item ก) โมเมนตัมคงที่
\item ข) ความเร็วคงที่
\item ค) การเปลี่ยนแปลงโมเมนตัม
\item ง) วัตถุเคลื่อนที่เป็นเส้นตรงด้วยความเร็วคงที่
\end{enumerate}
\item \textbf{เรือลำหนึ่งมวล 200 kg ลอยนิ่งบนน้ำ นักว่ายน้ำมวล 60 kg กระโดดออกจากเรือด้วยความเร็ว 4 m/s ความเร็วของเรือเป็นเท่าใด}
\begin{enumerate}
\item ก) 0.8 m/s ตามทิศที่นักว่ายกระโดด
\item ข) 1.0 m/s ตามทิศที่นักว่ายกระโดด
\item ค) 1.2 m/s ตรงข้ามทิศที่นักว่ายกระโดด
\item ง) 1.5 m/s ตรงข้ามทิศที่นักว่ายกระโดด
\end{enumerate}
\item \textbf{ลูกบอลมวล 0.4 kg กระทบผนังด้วยมุม 60° และสะท้อนออกด้วยมุม 60° ด้วยขนาดความเร็วเท่าเดิม ความเร็วก่อนชน = 10 m/s impulse มีค่าเท่าใด}
\begin{enumerate}
\item ก) 2 N·s
\item ข) 4 N·s
\item ค) 6 N·s
\item ง) 8 N·s
\end{enumerate}
\item \textbf{ลูกกอล์ฟมวล 0.05 kg วิ่งด้วยความเร็ว 40 m/s ชนเข้ากับไม้กอล์ฟแล้วอยู่ติดกับไม้ ถ้ามวลของไม้กอล์ฟ 0.95 kg และแรงเฉลี่ยที่ไม้กระทำ = 2000 N จะใช้เวลาชนนานเท่าใด}
\begin{enumerate}
\item ก) 0.0005 s
\item ข) 0.001 s
\item ค) 0.002 s
\item ง) 0.005 s
\end{enumerate}
\end{enumerate}
\end{document}